\section*{Bab \ref{ch:b3:clt} --- Fungsi
Karakteristik dan Teorema Limit Pusat}

\begin{solution}{exo:b3:clt:1}
Seragam pada $\intcc{-1}1$: $\varphi(\xi) =
\frac12\int_{-1}^1\eu^{\iu\xi x}\dd x =
\frac{\sin\xi}{\xi}$ (yang sama dengan $1$ di $\xi = 0$).
Eksponensial $\mathcal E(\lambda)$: $\varphi(\xi) =
\lambda\int_0^\infty\eu^{(\iu\xi - \lambda)x}\dd x =
\frac{\lambda}{\lambda - \iu\xi}$ (sebab antiturunannya
lenyap di $+\infty$ karena $\operatorname{Re}(\iu\xi -
\lambda) < 0$). Poisson $\mathcal P(\lambda)$: menurut
teorema transfer bagi distribusi diskretnya,
\[
\varphi(\xi) = \sum_{k\geq0}\eu^{\iu\xi
k}\,\eu^{-\lambda}\frac{\lambda^k}{k!}
= \eu^{-\lambda}\exp\bigl(\lambda\eu^{\iu\xi}\bigr)
= \exp\bigl(\lambda(\eu^{\iu\xi} - 1)\bigr).
\]
Binomial $\mathcal B(n, p)$: yakni jumlah $n$ variabel
Bernoulli yang bebas, masing-masing berfungsi karakteristik $1 - p + p\eu^{\iu\xi}$,
sehingga $\varphi(\xi) = \bigl(1 - p + p\eu^{\iu\xi}\bigr)^n$
(\cref{prop:b3:clt:cfbasics}(b)). Untuk keaditifan Poissonnya: jika $X
\sim \mathcal P(\lambda)$ dan $Y \sim \mathcal P(\mu)$
saling bebas,
\[
\varphi_{X+Y}(\xi) = \eu^{\lambda(\eu^{\iu\xi}-1)}
\eu^{\mu(\eu^{\iu\xi}-1)} =
\eu^{(\lambda+\mu)(\eu^{\iu\xi}-1)},
\]
yakni \omterm{def:b3:clt:cf}{fungsi karakteristik} $\mathcal P(\lambda + \mu)$; lalu keinjektifannya
(\cref{thm:b3:clt:injectivity}) mengenali distribusinya.
\end{solution}

\begin{solution}{exo:b3:clt:2}
(a) Berlaku $\overline{\varphi_X(\xi)} = \E\eu^{-\iu\xi X} =
\varphi_{-X}(\xi)$. Jadi $\varphi_X$ real jika dan hanya jika
$\varphi_X = \varphi_{-X}$, jika dan hanya jika (menurut keinjektifannya,
\cref{thm:b3:clt:injectivity}) $X$ dan $-X$ berdistribusi
sama.
(b) Tulislah $\varphi_X(\xi_0) = \eu^{\iu\theta}$. Maka
\[
\E\bigl[1 - \cos(\xi_0X - \theta)\bigr]
= 1 - \operatorname{Re}\bigl(\eu^{-\iu\theta}
\varphi_X(\xi_0)\bigr) = 1 - 1 = 0 .
\]
Integrannya tak negatif, sehingga $\cos(\xi_0X - \theta) = 1$
hampir pasti (sebab variabel tak negatif yang bernilai harapan nol
lenyap secara h.p.), yakni $\xi_0X - \theta \in 2\pi\Z$ secara h.p.: jadi $X$
mengambil nilainya pada barisan aritmetika
$\frac{\theta}{\xi_0} + \frac{2\pi}{\xi_0}\Z$ hampir pasti.
Sedangkan jika $X$ berkepadatan, himpunan terbilang ini bernilai Lebesgue nol,
sehingga berpeluang $0$ --- yakni kontradiksi; karena itu
$\abs{\varphi_X(\xi)} < 1$ bagi setiap $\xi \neq 0$.
\end{solution}

\begin{solution}{exo:b3:clt:3}
Menurut \omterm{def:b3:probability:independence}{kebebasannya} dan \cref{prop:b3:clt:cfbasics}:
\[
\varphi_{X+Y}(\xi) = \eu^{\iu m_1\xi - \sigma_1^2\xi^2/2}\,
\eu^{\iu m_2\xi - \sigma_2^2\xi^2/2}
= \eu^{\iu(m_1+m_2)\xi - (\sigma_1^2+\sigma_2^2)\xi^2/2},
\]
yakni \omterm{def:b3:clt:cf}{fungsi karakteristik} $\mathcal N(m_1 + m_2, \sigma_1^2 + \sigma_2^2)$;
lalu keinjektifannya menyimpulkan. Sedangkan kestabilan terhadap pemetaan afin adalah
aturan afinnya ($aX + b \sim \mathcal N(am_1 + b,
a^2\sigma_1^2)$, dengan mengizinkan kasus merosot $a = 0$), dan
kestabilan terhadap jumlah yang bebas menyusul lewat induksi atas
perhitungan di atas. Adapun bagi \omterm{def:b3:clt:gaussianvector}{Gauss} yang \emph{bergantung}, jumlahnya
tak harus \omterm{def:b3:clt:gaussianvector}{Gauss}: sebab pada \cref{exo:b3:clt:9}, $X$ dan $Y =
\varepsilon X$ masing-masing \omterm{def:b3:clt:gaussianvector}{Gauss} baku tetapi $X + Y$
lenyap dengan peluang $\frac12$ tanpa nol secara
h.p., sehingga ia bukan \omterm{def:b3:clt:gaussianvector}{Gauss}.
\end{solution}

\begin{solution}{exo:b3:clt:4}
(a) \emph{Implikasi langsungnya.} Misalkan $t$ titik \omterm{def:b3:topology:continuity}{kekontinuan}
$F_X$ dan $\delta > 0$. Ambillah landaian \omterm{def:b3:topology:continuity}{kontinu}
$f^-$ (yang $= 1$ pada $\intoc{-\infty}{t-\delta}$, $0$ mulai dari $t$,
dan afin di antaranya) dan $f^+$ (yang $= 1$ pada $\intoc{-\infty}t$,
$0$ mulai dari $t + \delta$, dan afin di antaranya); maka $f^- \leq
\mathbf 1_{\intoc{-\infty}t} \leq f^+$, sehingga
\[
\E f^-(X_n) \leq F_{X_n}(t) \leq \E f^+(X_n),
\]
sedangkan suku luarnya konvergen ke $\E f^\pm(X)$, yang sendirinya
terapit di antara $F_X(t - \delta)$ dan $F_X(t + \delta)$.
Lalu dengan membiarkan $n \to \infty$ kemudian $\delta \to 0$ dan memakai
\omterm{def:b3:topology:continuity}{kekontinuan} $F_X$ di $t$: $F_{X_n}(t) \to F_X(t)$.

\emph{Konversnya.} Misalkan $f$ \omterm{def:b3:topology:continuity}{kontinu} terbatas, $M =
\sup\abs f$, dan $\varepsilon > 0$. Titik \omterm{def:b3:topology:continuity}{kekontinuan}
$F_X$ bersifat padat (sebab $F_X$ berlompatan paling banyak terbilang), sehingga
pilihlah titik \omterm{def:b3:topology:continuity}{kekontinuan} $a < b$ dengan $F_X(a) < \varepsilon$
dan $1 - F_X(b) < \varepsilon$. Lalu pada $\intcc ab$ yang \omterm{def:b3:topology:compact}{kompak},
fungsi $f$ \omterm{def:b3:topology:continuity}{kontinu} seragam: pilihlah titik
\omterm{def:b3:topology:continuity}{kekontinuan} $a = t_0 < t_1 < \dots < t_m = b$ milik $F_X$ yang
ayunan $f$-nya paling banyak $\varepsilon$ pada setiap
$\intoc{t_{j-1}}{t_j}$, lalu tetapkan $g = \sum_j
f(t_j)\,\mathbf 1_{\intoc{t_{j-1}}{t_j}}$. Maka $\abs{f - g}
\leq \varepsilon$ pada $\intoc ab$, $\abs g \leq M$, dan bagi
$T = X_n$ atau $X$:
\[
\bigl|\E f(T) - \E g(T)\bigr| \leq \varepsilon +
2M\bigl(F_T(a) + 1 - F_T(b)\bigr).
\]
Lebih jauh $\E g(X_n) = \sum_j f(t_j)\bigl(F_{X_n}(t_j) -
F_{X_n}(t_{j-1})\bigr) \to \E g(X)$ (yakni jumlah berhingga suku
yang konvergen, sebab semua $t_j$ merupakan titik \omterm{def:b3:topology:continuity}{kekontinuannya}), dan
$F_{X_n}(a) \to F_X(a) < \varepsilon$, $1 - F_{X_n}(b) \to 1
- F_X(b) < \varepsilon$. Lalu dengan merakitnya:
$\limsup_n\abs{\E f(X_n) - \E f(X)} \leq 2\varepsilon +
8M\varepsilon$; lalu biarkanlah $\varepsilon \to 0$.

(b) Fungsi distribusi konstanta $c$ adalah
$\mathbf 1_{\intco c\infty}$, yang \omterm{def:b3:topology:continuity}{kontinu} kecuali di $c$. Lalu bagi
$\varepsilon > 0$, titik $c - \varepsilon$ dan $c +
\frac\varepsilon2$ merupakan titik \omterm{def:b3:topology:continuity}{kekontinuannya}, sehingga
\[
\P(\abs{X_n - c} > \varepsilon) \leq F_{X_n}(c -
\varepsilon) + 1 - F_{X_n}\Bigl(c + \frac\varepsilon2\Bigr)
\longrightarrow 0 + 1 - 1 = 0 .
\]
\end{solution}

\begin{solution}{exo:b3:clt:5}
Ambillah $z_n = p_n(\eu^{\iu\xi} - 1)$, sehingga $\varphi_{X_n}(\xi) =
(1 + z_n)^n$ (\cref{exo:b3:clt:1}) dan $\abs{z_n} \leq 2p_n
\to 0$ (perhatikan $p_n = \frac{np_n}n \to 0$). Baik $1 + z_n$ maupun
$\eu^{z_n}$ bermodulus paling banyak $1$: sebab $\abs{1 + z_n} =
\abs{(1 - p_n) + p_n\eu^{\iu\xi}} \leq 1$ menurut ketaksamaan
segitiga, sedangkan $\abs{\eu^{z_n}} = \eu^{p_n(\cos\xi - 1)}
\leq 1$. Lalu ketaksamaan teleskopnya $\abs{a^n - b^n} \leq
n\abs{a - b}$ (yakni bukti \cref{thm:b3:clt:clt}) dan batas
deret pangkatnya $\abs{\eu^z - 1 - z} \leq
\abs z^2\eu^{\abs z}$ memberikan
\[
\bigl|(1 + z_n)^n - \eu^{nz_n}\bigr| \leq n\bigl|1 + z_n -
\eu^{z_n}\bigr| \leq n\,\abs{z_n}^2\,\eu^{\abs{z_n}} \leq
4\eu^2\,np_n^2 = 4\eu^2\,(np_n)\,p_n \longrightarrow 0 .
\]
Karena $nz_n = np_n(\eu^{\iu\xi} - 1) \to
\lambda(\eu^{\iu\xi} - 1)$, kita simpulkan $\varphi_{X_n}(\xi)
\to \exp\bigl(\lambda(\eu^{\iu\xi} - 1)\bigr)$ bagi setiap
$\xi$: yakni \omterm{def:b3:clt:cf}{fungsi karakteristik} $\mathcal P(\lambda)$, sehingga Lévy
(\cref{thm:b3:clt:levy}) memberikan $X_n \Rightarrow \mathcal
P(\lambda)$. Secara numerik: $\P\bigl(\mathcal B(100, 0.02) =
0\bigr) = 0.98^{100} = \eu^{100\ln 0.98} \approx
\eu^{-2.020} \approx 0.1326$, sedangkan $\P\bigl(\mathcal P(2) =
0\bigr) = \eu^{-2} \approx 0.1353$: jadi selisihnya dua persen
sudah pada $n$ yang kasar ini.
\end{solution}

\begin{solution}{exo:b3:clt:6}
(a) Satu lemparan berrata-rata $\frac72$ dan bervarians
$\frac{35}{12}$, sehingga $S$ berrata-rata $3500$, bervarians
$\frac{35000}{12} \approx 2916.7$ dan bersimpangan baku
$\approx 54.0$. Lalu menurut teorema limit pusatnya,
\[
\P(S > 3600) = \P\Bigl(\frac{S - 3500}{54.0} > 1.85\Bigr)
\approx 1 - \Phi(1.85) \approx 0.032 :
\]
yakni peluang sekitar $3\%$.
(b) Untuk $S \sim \mathcal B(100, \frac12)$: rata-ratanya $50$ dan simpangan
bakunya $5$. Dengan koreksi \omterm{def:b3:topology:continuity}{kekontinuannya},
\[
\P(45 \leq S \leq 55) \approx
\Phi\Bigl(\frac{55.5 - 50}{5}\Bigr) -
\Phi\Bigl(\frac{44.5 - 50}{5}\Bigr) = 2\Phi(1.1) - 1
\approx 0.729,
\]
terhadap nilai persisnya $0.7287$; sedangkan tanpa koreksinya,
$2\Phi(1) - 1 \approx 0.683$, yang meleset hampir lima poin. Adapun
koreksinya penting karena $S$ merupakan variabel kisi: sebab
atom $\P(S = k)$ terhampiri baik oleh massa \omterm{def:b3:clt:gaussianvector}{Gauss}
$\intcc{k - \frac12}{k + \frac12}$, sehingga menggunting
selangnya di bilangan bulat $45$ dan $55$ membuang separuh atom
pada setiap ujungnya.
\end{solution}

\begin{solution}{exo:b3:clt:7}
Variabel $g(U_k)$ bersifat i.i.d.\ (yakni peta \omterm{def:b3:lebesgue:measurable}{terukur}
variabel i.i.d.), \omterm{def:b3:lebesgue:l1}{terintegralkan} kuadrat, berrata-rata $I$
(menurut teorema transfer, \cref{exo:b3:product:9}) dan bervarians
$\sigma^2$. Jika $\sigma > 0$, maka \cref{thm:b3:clt:clt} yang diterapkan
padanya persis merupakan konvergensi yang dinyatakan
\[
\sqrt n\,\Bigl(\frac1n\sum_{k\leq n}g(U_k) - I\Bigr) =
\frac{\sum_{k\leq n}\bigl(g(U_k) - I\bigr)}{\sqrt n}
\Longrightarrow \mathcal N(0, \sigma^2)
\]
(sedangkan jika $\sigma = 0$, maka $g$ konstan secara h.p.\ dan ruas
kirinya lenyap secara identik). Karena itu
$\P\bigl(\abs{\frac1n\sum g(U_k) - I} \leq
1.96\,\sigma/\sqrt n\bigr) \to 0.95$: jadi batang galat $\pm
1.96\,\sigma/\sqrt n$ melihat dimensi $d$ hanya lewat
konstanta $\sigma$-nya, tak pernah lewat lajunya terhadap $n$. Sedangkan
kaidah titik tengah dengan $n$ simpul pada dimensi $d$ berjala
$n^{-1/d}$ dan bergalat berorde $n^{-2/d}$ bagi integran
$\mathcal C^2$. Lalu $n^{-1/2}$ milik \omterm{ex:b3:probability:sllnapps}{Monte Carlo} meluruh lebih cepat daripada
$n^{-2/d}$ tepat saat $\frac12 > \frac2d$, yakni $d > 4$:
jadi mulai dimensi $5$, pencuplikan acaknya secara asimtotik mengalahkan
kisinya --- sebab kutukan dimensinya mengampuni metode
peluang, dan itulah sebabnya \omterm{ex:b3:probability:sllnapps}{Monte Carlo} menguasai
pengintegralan berdimensi tinggi.
\end{solution}

\begin{solution}{exo:b3:clt:8}
\emph{Jumlahnya.} Bagi $\xi$ yang tetap:
\[
\bigl|\E\eu^{\iu\xi(X_n+Y_n)} -
\eu^{\iu\xi c}\,\E\eu^{\iu\xi X_n}\bigr|
= \bigl|\E\bigl[\eu^{\iu\xi X_n}\bigl(\eu^{\iu\xi Y_n} -
\eu^{\iu\xi c}\bigr)\bigr]\bigr|
\leq \E\bigl|\eu^{\iu\xi(Y_n - c)} - 1\bigr| .
\]
Belahlah pada kejadian $\{\abs{Y_n - c} \leq \delta\}$: di sana
$\abs{\eu^{\iu\xi(Y_n-c)} - 1} \leq \abs\xi\,\delta$ (sebab
talinya lebih pendek daripada busurnya); sedangkan komplemennya menyumbang
paling banyak $2\,\P(\abs{Y_n - c} > \delta) \to 0$. Karena itu
$\limsup$-nya $\leq \abs\xi\,\delta$ bagi setiap $\delta > 0$:
jadi selisihnya menuju $0$. Lalu karena $\E\eu^{\iu\xi X_n} \to
\varphi_X(\xi)$, kita peroleh $\varphi_{X_n+Y_n}(\xi) \to
\eu^{\iu\xi c}\varphi_X(\xi) = \varphi_{X+c}(\xi)$, sehingga
Lévy (\cref{thm:b3:clt:levy}) menghasilkan $X_n + Y_n
\Rightarrow X + c$.

\emph{Hasil kalinya.} Pertama, $cX_n \Rightarrow cX$: sebab
$\varphi_{cX_n}(\xi) = \varphi_{X_n}(c\xi) \to
\varphi_X(c\xi) = \varphi_{cX}(\xi)$. Berikutnya, $(Y_n - c)X_n
\to 0$ dalam peluang: sebab distribusi $X_n$ bersifat ketat (karena
\omterm{def:b3:clt:cf}{fungsi karakteristiknya} konvergen ke sebuah \omterm{def:b3:clt:cf}{fungsi karakteristik}; lihat langkah keketatan
\cref{thm:b3:clt:levy}), sehingga diberikan $\varepsilon > 0$ pilihlah $M$
dengan $\P(\abs{X_n} > M) \leq \varepsilon$ bagi setiap $n$; maka
\[
\P\bigl(\abs{(Y_n - c)X_n} > \varepsilon\bigr) \leq
\P(\abs{X_n} > M) + \P\Bigl(\abs{Y_n - c} >
\frac{\varepsilon}{M}\Bigr) \leq \varepsilon + o(1) .
\]
Lalu dengan menulis $Y_nX_n = cX_n + (Y_n - c)X_n$ dan menerapkan
bagian jumlahnya (yang buktinya hanya memakai $Y_n' := (Y_n - c)X_n \to 0$ dalam
peluang, dengan konstanta $0$): jadi $Y_nX_n \Rightarrow cX$.

\emph{Penerapannya.} Menurut hukum kuat bilangan besar
(\cref{thm:b3:probability:slln}), $\hat p_n \to p$ secara h.p., sehingga
menurut \omterm{def:b3:topology:continuity}{kekontinuannya} $\hat\sigma_n = \sqrt{\hat p_n(1 - \hat p_n)}
\to \sigma = \sqrt{p(1 - p)} > 0$ secara h.p., jadi
$\frac{\sigma}{\hat\sigma_n} \to 1$ dalam peluang. Lalu aturan hasil kali
Slutsky menaikkan $\frac{S_n - np}{\sigma\sqrt n}
\Rightarrow \mathcal N(0,1)$ menjadi $\frac{S_n -
np}{\hat\sigma_n\sqrt n} = \frac{\sigma}{\hat\sigma_n}\cdot
\frac{S_n - np}{\sigma\sqrt n} \Rightarrow \mathcal N(0,1)$:
yakni selang kepercayaan yang \emph{terpakai} $\hat p_n \pm
1.96\,\hat\sigma_n/\sqrt n$, yang dibangun dari datanya saja,
dan tetap beraras asimtotik $95\%$.
\end{solution}

\begin{solution}{exo:b3:clt:9}
(a) Dengan membelah \omterm{def:b3:probability:space}{nilai harapannya} atas kedua nilai
$\varepsilon$ (menurut \omterm{def:b3:probability:independence}{kebebasannya}): bagi $B$ yang Borel, $\P(Y \in B) =
\frac12\P(X \in B) + \frac12\P(-X \in B) = \P(X \in B)$,
sebab $-X \sim X$ (karena $\mathcal N(0,1)$ setangkup): jadi $Y \sim
\mathcal N(0,1)$. Dan $\operatorname{Cov}(X, Y) =
\E[\varepsilon X^2] = \E[\varepsilon]\,\E[X^2] = 0 \cdot 1 =
0$.
(b) Berlaku $\abs Y = \abs X$, sehingga $\P(\abs X \leq 1,\ \abs Y \geq 2)
= 0$ padahal $\P(\abs X \leq 1)\,\P(\abs Y \geq 2) > 0$: jadi tak
bebas. Seandainya $(X, Y)$ merupakan \omterm{def:b3:clt:gaussianvector}{vektor Gauss}, maka $X + Y =
(1 + \varepsilon)X$ akan menjadi variabel \omterm{def:b3:clt:gaussianvector}{Gauss} real
(\cref{def:b3:clt:gaussianvector} dengan $t = (1,1)$); padahal
$\P(X + Y = 0) = \P(\varepsilon = -1) = \frac12$, sedangkan variabel
\omterm{def:b3:clt:gaussianvector}{Gauss} beratom hanya jika ia konstan secara h.p.\
--- dan $X + Y$ sama dengan $2X \neq 0$ secara h.p.\ pada
$\{\varepsilon = 1\}$. Jadi bertentangan: sehingga $(X, Y)$ bukan
\omterm{def:b3:clt:gaussianvector}{Gauss}.
(c) Setiap marginalnya \omterm{def:b3:clt:gaussianvector}{Gauss} dan kovariansnya lenyap,
padahal \omterm{def:b3:probability:independence}{kebebasannya} gagal --- sebab \emph{pasangannya} tak
\omterm{def:b3:clt:gaussianvector}{Gauss} bersama. Jadi \cref{thm:b3:clt:gaussianvector}(2) tak dapat
dilemahkan menjadi ``marginal \omterm{def:b3:clt:gaussianvector}{Gauss}''.
\end{solution}

\begin{solution}{exo:b3:clt:10}
(a) \cref{exo:b3:fouriertransform:1} menghitung
$\widehat{\eu^{-\abs\cdot}}(\xi) = \frac{2}{1 + \xi^2}$;
lalu karena kedua ruasnya \omterm{def:b3:lebesgue:l1}{terintegralkan}, pembalikan Fourier
(\cref{thm:b3:fouriertransform:inversion}) membalikkan
pernyataannya:
\[
\int_\R\eu^{\iu\xi x}\,\frac{\dd x}{\pi(1 + x^2)} =
\eu^{-\abs\xi},
\]
yang persis merupakan $\varphi_X(\xi)$ bagi variabel Cauchy $X$.
(b) Menurut \omterm{def:b3:probability:independence}{kebebasannya}, $\varphi_{S_n}(\xi) =
\bigl(\eu^{-\abs\xi}\bigr)^n = \eu^{-n\abs\xi}$, sehingga
$\varphi_{S_n/n}(\xi) = \varphi_{S_n}(\xi/n) =
\eu^{-\abs\xi}$: jadi rata-rata empiris $\frac{S_n}n$ kembali
Cauchy baku bagi setiap $n$ (menurut keinjektifannya). Sehingga rata-ratanya
tak pernah memusat: sebab fluktuasinya pada waktu $10^6$ sama dengan
fluktuasi satu pengamatan.
(c) Berlaku $\E\abs{X_1} = \frac2\pi\int_0^\infty\frac{x\,\dd x}{1 +
x^2} = +\infty$: jadi distribusi Cauchy tak \omterm{def:b3:lebesgue:l1}{terintegralkan}, sehingga
hukum kuat bilangan besar
(\cref{thm:b3:probability:slln}) tak berlaku, dan teorema limit pusatnya
(yang menuntut varians berhingga) apalagi. Jadi di sini
kesimpulannya sungguh gagal, bukan sekadar buktinya.
Sebagai uji keselarasan: $\varphi(\xi) = \eu^{-\abs\xi}$ tak
terdiferensialkan di $0$, seperti yang diramalkan \cref{prop:b3:clt:cfbasics}(c)
yang dibaca secara kontraposisi bagi variabel yang tak
\omterm{def:b3:lebesgue:l1}{terintegralkan}.
\end{solution}

\begin{solution}{exo:b3:clt:11}
(a) Berlaku $\varphi_{aX_1 + bX_2}(\xi) =
\eu^{-a\abs\xi}\eu^{-b\abs\xi} = \eu^{-(a+b)\abs\xi} =
\varphi_{(a+b)X_1}(\xi)$ (menurut \omterm{def:b3:probability:independence}{kebebasannya} dan
\cref{exo:b3:clt:10}); lalu keinjektifannya mengenali distribusinya.

(b) Berlaku $\varphi_{aX_1+bX_2}(\xi) = \eu^{-a^2\xi^2/2}
\eu^{-b^2\xi^2/2} = \eu^{-(a^2+b^2)\xi^2/2}$: yakni distribusi
$\sqrt{a^2+b^2}\,X_1$.

(c) Jika $\varphi_X(\xi) = \eu^{-c\abs\xi^\alpha}$, maka
$S_n = X_1 + \dots + X_n$ mempunyai $\varphi_{S_n} =
\eu^{-cn\abs\xi^\alpha}$, sedangkan $S_n/n^{1/\alpha}$ mempunyai
$\varphi(\xi) = \eu^{-c\abs\xi^\alpha}$ lagi: yakni
peranakan diri yang persis di bawah penskalaan $n^{1/\alpha}$ ---
$\sqrt n$ bagi Gaussnya ($\alpha = 2$), dan $n$ itu sendiri bagi
Cauchynya ($\alpha = 1$, \cref{exo:b3:clt:10}(b)). Jadi jumlah
variabel i.i.d.\ hanya dapat konvergen (setelah penormalan
afin) ke distribusi yang stabil terhadap konvolusi
semacam itu; sedangkan teorema limit pusatnya mengatakan bahwa varians berhingga memaksa
cekungan Gaussnya, dan cekungan Cauchynya tersedia bagi distribusi
yang ekornya begitu berat sehingga $\E X^2 = \infty$ bahkan $\E\abs
X = \infty$ --- misalnya jumlah variabel Cauchy itu sendiri.
Jadi keuniversalannya berpulau-pulau, yang terindeks eksponen
ekornya $\alpha \in \intoc02$.
\end{solution}

\begin{solution}{exo:b3:clt:12}
(a) Indikator $\mathbf 1_{X_k \leq t}$ bersifat i.i.d.\
Bernoulli berparameter $p = F(t)$: sehingga jumlahnya $nF_n(t)$
binomial $\mathcal B(n, p)$; lalu hukum kuatnya memberikan $F_n(t)
\to p$ secara h.p., sedangkan teorema limit pusatnya (\cref{thm:b3:clt:clt}) yang diterapkan pada
indikator yang sama (bervarians $p(1-p)$) memberikan limit
\omterm{def:b3:clt:gaussianvector}{Gauss} yang dinyatakan itu.

(b) Nilai $p(1 - p)$ maksimal di $p = \frac12$, yakni di tempat
$F(t) = \frac12$: yaitu di \emph{mediannya}. Jadi menaksir peluang
ekornya mudah secara asimtotik (sebab variansnya $\to 0$ saat
$p \to 0, 1$); sedangkan daerah mediannya mengangkut derau
statistik terbesar --- yakni kurva empirisnya paling bergoyang di
tengahnya.

(c) Untuk $F$ yang \omterm{def:b3:topology:continuity}{kontinu}, variabel $U_k = F(X_k)$ bersifat
seragam i.i.d.\ pada $\intoo01$ (\cref{exo:b3:probability:1}),
lalu kemonotonan $F$ memberikan, dengan menulis $G_n$ bagi
fungsi distribusi empiris $U_k$-nya:
\[
\sup_{t\in\R}\,\abs{F_n(t) - F(t)}
= \sup_{u \in \operatorname{im}F}\,\abs{G_n(u) - u}
= \sup_{u\in\intcc01}\abs{G_n(u) - u} :
\]
yakni kesamaan pertamanya karena $\{X_k \leq t\} = \{U_k \leq
F(t)\}$ sampai kejadian nol (menurut kemonotonannya; sebab ketaksamaan sejatinya
hanya dapat gagal pada bagian datar $F$, yang kedua ruasnya
tak berubah di sana), sedangkan yang kedua karena $F$ yang \omterm{def:b3:topology:continuity}{kontinu},
yang berjalan dari $0$ ke $1$, mencapai setiap nilai
$\intoo01$ (menurut teorema nilai antaranya), dan ujungnya
tak menambah apa pun (sebab $G_n(0) - 0 = 0$ dan $G_n(1) - 1 = 0$). Jadi
ruas kanannya hanya melibatkan variabel seragam: yakni satu distribusi
bagi setiap $F$ --- sehingga satu tabel nilai kritis (yakni tabel
distribusi Kolmogorov) menguji \emph{sembarang} model \omterm{def:b3:topology:continuity}{kontinu}
terhadap datanya.
\end{solution}

\begin{solution}{pb:b3:clt:1}
\textbf{1.} Keluarga $(X_1, \dots, X_n, N_1, \dots, N_n)$
bersifat bebas: sebab kedua bloknya saling bebas
menurut konstruksinya dan setiap bloknya i.i.d. Lalu $W_i$ merupakan
\omterm{def:b3:lebesgue:measurable}{fungsi terukur} variabel $(X_j)_{j<i}$ dan
$(N_j)_{j>i}$ saja, yang semuanya berbeda dari $X_i$ dan $N_i$: sehingga menurut
asas koalisinya
(\cref{thm:b3:probability:independence}), $W_i$ bersifat
bebas dari pasangan $(X_i, N_i)$. Adapun penguraian
$H_i = W_i + \frac{X_i}{\sqrt n}$ dan $H_{i-1} = W_i +
\frac{N_i}{\sqrt n}$ langsung dari definisinya:
sebab berpindah dari $H_i$ ke $H_{i-1}$ menukar satu suku
$X_i$ dengan $N_i$.

\textbf{2.} Taylor--Lagrange pada orde $3$: ada $c$
di antara $w$ dan $w + h$ dengan $f(w + h) = f(w) + f'(w)h +
\frac12f''(w)h^2 + \frac16f'''(c)h^3$, lalu $\abs{f'''(c)}
\leq M_3$ memberikan batasnya.

\textbf{3.} Dengan mengurangkan kedua uraiannya pada titik
pangkal bersamanya $w = W_i$:
\[
f(H_i) - f(H_{i-1}) = f'(W_i)\,\frac{X_i - N_i}{\sqrt n} +
\frac{f''(W_i)}{2}\,\frac{X_i^2 - N_i^2}{n} + R_i,
\qquad
\abs{R_i} \leq \frac{M_3}{6}\cdot
\frac{\abs{X_i}^3 + \abs{N_i}^3}{n^{3/2}} .
\]
Lalu ambillah \omterm{def:b3:probability:space}{nilai harapannya}. Menurut pertanyaan 1, $f'(W_i)$ dan $f''(W_i)$
bebas dari $(X_i, N_i)$, sehingga \omterm{def:b3:probability:space}{nilai harapan} campurannya
memfaktor:
\begin{align*}
\E\Bigl[f'(W_i)\,\frac{X_i - N_i}{\sqrt n}\Bigr] &=
\E\bigl[f'(W_i)\bigr]\,\frac{\E X_i - \E N_i}{\sqrt n} = 0,
\\
\E\Bigl[f''(W_i)\,\frac{X_i^2 - N_i^2}{n}\Bigr] &=
\E\bigl[f''(W_i)\bigr]\,\frac{1 - 1}{n} = 0 :
\end{align*}
yakni dua momen pertama $X_i$ dan $N_i$ \emph{cocok}, sehingga
hanya sisanya yang bertahan:
\[
\bigl|\E f(H_i) - \E f(H_{i-1})\bigr| \leq \E\abs{R_i} \leq
\frac{M_3}{6}\cdot\frac{\beta + \gamma}{n^{3/2}} .
\]
Adapun momen ketiga Gaussnya: $\gamma = \E\abs{N_1}^3 =
2\int_0^\infty x^3\,\frac{\eu^{-x^2/2}}{\sqrt{2\pi}}\,\dd x
= \frac{2}{\sqrt{2\pi}}\int_0^\infty 2u\,\eu^{-u}\dd u =
\frac{4}{\sqrt{2\pi}} = \frac{2\sqrt2}{\sqrt\pi}$
(lewat substitusi $u = x^2/2$, lalu $\Gamma(2) = 1$).

\textbf{4.} Dengan menteleskopkan $\E f(T_n) - \E f(G_n) =
\sum_{i=1}^n\bigl(\E f(H_i) - \E f(H_{i-1})\bigr)$ lalu
menerapkan pertanyaan 3 pada masing-masing $n$ sukunya:
\[
\bigl|\E f(T_n) - \E f(G_n)\bigr| \leq
n \cdot \frac{M_3(\beta + \gamma)}{6\,n^{3/2}} =
\frac{M_3\,(\beta + \gamma)}{6\,\sqrt n} .
\]

\textbf{5.} Variabel $G_n$ bersifat $\mathcal N(0,1)$ \emph{secara persis} bagi
setiap $n$ (sebab jumlah ternormalkan variabel \omterm{def:b3:clt:gaussianvector}{Gauss} baku yang
bebas, \cref{exo:b3:clt:3}), sehingga $\E f(G_n) = \E f(N)$
dan pertanyaan 4 berbunyi $\abs{\E f(T_n) - \E f(N)} \leq
\frac{M_3(\beta+\gamma)}{6\sqrt n} \to 0$ bagi $f \in
\mathcal C^3_b$. \emph{Untuk kenaikannya.} Misalkan $f$ \omterm{def:b3:topology:continuity}{kontinu}
terbatas, $M = \sup\abs f$, dan $\varepsilon > 0$. Pilihlah $A
\geq 1$ dengan $\frac1{A^2} \leq \varepsilon$: lalu Chebyshev beserta
$\V(T_n) = 1$ memberikan $\P(\abs{T_n} > A) \leq \varepsilon$
bagi setiap $n$, dan demikian pula $\P(\abs N > A) \leq \varepsilon$.
Misalkan $\chi$ bersifat $\mathcal C^\infty$ dengan
$\mathbf 1_{\intcc{-A}A} \leq \chi \leq
\mathbf 1_{\intcc{-A-1}{A+1}}$ (yakni dataran mulus, yang dibangun dengan
memuluskan $\mathbf 1_{\intcc{-A-\frac12}{A+\frac12}}$,
\cref{thm:b3:lp:regularization}); maka $g = f\chi$ \omterm{def:b3:topology:continuity}{kontinu}
bertumpuan \omterm{def:b3:topology:compact}{kompak}, jadi \omterm{def:b3:topology:continuity}{kontinu} seragam, sehingga
pemulusannya $g_\eta = g * \rho_\eta$ bersifat $\mathcal
C^\infty$ dengan turunan terbatas pada setiap orde dan
$\norm{g - g_\eta}_\infty \leq \varepsilon$ bagi $\eta$ yang
cukup kecil. Lalu bagi $T = T_n$ atau $N$, karena $f = g$ pada
$\intcc{-A}A$ dan $\abs{f - g} \leq 2M$ di mana-mana:
\[
\bigl|\E f(T) - \E g_\eta(T)\bigr| \leq
\E\abs{(f - g)(T)} + \norm{g - g_\eta}_\infty
\leq 2M\,\P(\abs T > A) + \varepsilon \leq
(2M + 1)\,\varepsilon .
\]
Lalu dengan menggabungkannya dengan $\E g_\eta(T_n) \to \E g_\eta(N)$ (sebab pertanyaan
4 berlaku: $g_\eta \in \mathcal C^3_b$):
\[
\limsup_n\;\bigl|\E f(T_n) - \E f(N)\bigr| \leq
2(2M + 1)\,\varepsilon ,
\]
sedangkan $\varepsilon$-nya sembarang: jadi $\E f(T_n) \to \E f(N)$ bagi
setiap $f$ yang \omterm{def:b3:topology:continuity}{kontinu} terbatas, yakni $T_n \Rightarrow
\mathcal N(0,1)$.

\textbf{6.} Distribusi yang identik masuk hanya lewat satu
kalimat: ``$X_i$ dan $N_i$ mempunyai dua momen pertama yang
sama''. Jadi misalkan $X_1, \dots, X_n$ bebas,
terpusat, bervarians $\sigma_i^2$ dan bermomen ketiga
berhingga, dengan $s_n^2 = \sum_i\sigma_i^2 > 0$, lalu ambillah $N_i \sim
\mathcal N(0, \sigma_i^2)$ yang bebas dari segalanya.
Definisikanlah hibridanya dengan penormalan $s_n$: $H_i =
\frac1{s_n}(\sum_{j\leq i}X_j + \sum_{j>i}N_j)$. Lalu pada
penukaran ke-$i$, $\E X_i = \E N_i = 0$ dan $\E X_i^2 = \E N_i^2
= \sigma_i^2$ kembali membunuh suku $f'$ dan $f''$-nya, sedangkan
sisanya memberikan (dengan memakai $\E\abs{N_i}^3 = \sigma_i^3\gamma$ lewat
penskalaan):
\[
\bigl|\E f(H_i) - \E f(H_{i-1})\bigr| \leq
\frac{M_3}{6\,s_n^3}\bigl(\E\abs{X_i}^3 +
\sigma_i^3\gamma\bigr) .
\]
Lalu dengan menteleskopkannya:
\[
\Bigl|\E f\Bigl(\frac{X_1 + \dots + X_n}{s_n}\Bigr) -
\E f(N)\Bigr| \leq \frac{M_3}{6\,s_n^3}\sum_{i=1}^n
\Bigl(\E\abs{X_i}^3 + \V(X_i)^{3/2}\,\gamma\Bigr) .
\]
Karena $\sigma_i^3 = (\E X_i^2)^{3/2} \leq \E\abs{X_i}^3$
(menurut ketaksamaan rata-rata pangkatnya, yakni Jensen bagi $t \mapsto
t^{3/2}$ yang diterapkan pada $X_i^2$), ruas kanannya paling banyak
$\frac{M_3(1 + \gamma)}{6}\cdot
\frac{\sum_i\E\abs{X_i}^3}{s_n^3}$: sehingga di bawah \emph{syarat
Lyapunov} $\frac1{s_n^3}\sum_i\E\abs{X_i}^3 \to 0$, jumlah
ternormalkannya konvergen dalam distribusi ke $\mathcal N(0,1)$ ---
yakni teorema limit pusatnya tanpa distribusi yang identik.

\textbf{7.} Misalkan $\rho \in \mathcal C^\infty_c(\intoo01)$
dengan $\int\rho = 1$ lalu tetapkan $\psi(x) = \int_x^1\rho(s)\dd
s$: maka $\psi$ bersifat $\mathcal C^\infty$, tak naik, dengan $\psi = 1$
pada $\R_-$ dan $\psi = 0$ pada $\intco1\infty$; lalu ambillah $K =
\norm{\psi'''}_\infty$. Lalu bagi $t \in \R$ dan $\delta > 0$
definisikanlah $\psi_\delta(x) = \psi\bigl(\frac{x - t}\delta\bigr)$
dan $\tilde\psi_\delta(x) = \psi\bigl(\frac{x - t}\delta +
1\bigr)$: keduanya $\mathcal C^3_b$ dengan turunan ketiga
terbatas oleh $K/\delta^3$, dan
\[
\mathbf 1_{\intoc{-\infty}{t-\delta}} \leq
\tilde\psi_\delta \leq \mathbf 1_{\intoc{-\infty}t} \leq
\psi_\delta \leq \mathbf 1_{\intoc{-\infty}{t+\delta}} .
\]
Untuk batas atasnya: menurut pertanyaan 4 yang diterapkan pada $\psi_\delta$ (dengan
$M_3 = K/\delta^3$),
\[
\P(T_n \leq t) \leq \E\psi_\delta(T_n) \leq
\E\psi_\delta(N) + \frac{K(\beta +
\gamma)}{6\,\delta^3\sqrt n}
\leq \Phi(t + \delta) + \frac{K(\beta +
\gamma)}{6\,\delta^3\sqrt n}
\leq \Phi(t) + \frac{\delta}{\sqrt{2\pi}} +
\frac{K(\beta + \gamma)}{6\,\delta^3\sqrt n},
\]
sebab $\Phi$ bersifat Lipschitz berkonstanta
$\frac1{\sqrt{2\pi}}$ (karena \omterm{ex:b3:lebesgue:gamma}{kepadatannya} terbatas oleh
$\frac1{\sqrt{2\pi}}$). Lalu batas bawah setangkupnya lewat
$\tilde\psi_\delta$ memberikan taksiran dua sukunya
\[
\sup_{t\in\R}\,\bigl|\P(T_n \leq t) - \Phi(t)\bigr| \leq
\frac{K(\beta + \gamma)}{6}\cdot\frac{1}{\delta^3\sqrt n} +
\frac{\delta}{\sqrt{2\pi}}
\qquad(\delta > 0\ \text{sembarang}).
\]
Kedua sukunya berimbang saat $\delta^{-3}n^{-1/2} \asymp
\delta$, yakni $\delta = n^{-1/8}$: sehingga keduanya menjadi
$O(n^{-1/8})$, yakni laju seragam eksplisit yang sahih bagi setiap
$n$. (Sedangkan laju Berry--Esseen yang optimum $C\beta/\sqrt n$
menuntut ketaksamaan pemulusan Fouriernya; jadi penukarannya menukar
ketajaman dengan keelementeran yang tuntas.)

\textbf{8.} Untuk $X_i = 2B_i - 1$ (yakni tanda adil): ia terpusat,
bervarians $1$, dan $\abs{X_i} = 1$ sehingga $\beta = 1$. Lalu pertanyaan 7
membatasi $\sup_t\abs{\P(\frac{S_n}{\sqrt n} \leq t) -
\Phi(t)}$ secara eksplisit dan seragam bagi \emph{setiap} $n$ yang
berhingga --- yakni pernyataan global dan tak asimtotik tentang
fungsi distribusinya. Sedangkan taksiran lokal
\cref{pb:b3:product:1}, pertanyaan 7, justru memberikan asimtotik
persis sebuah atom tunggal, $\P(S_{2n} = 2k) \sim
\frac{\eu^{-k^2/n}}{\sqrt{\pi n}}$: sehingga ia menguraikan
peluang berukuran $n^{-1/2}$, jauh di bawah resolusi $n^{-1/8}$
pertanyaan 7, tetapi ia bersifat titik demi titik, asimtotik
(tanpa galat eksplisit pada $n$ yang tetap) dan terikat pada
distribusi kisi tertentu ini. Jadi ketelitian lokal lawan
keseragaman global: kedua metodenya saling melengkapi, dan menjumlahkan
taksiran lokalnya atas $k \in \intint{a\sqrt n}{b\sqrt n}$
memulihkan de Moivre--Laplace pada selangnya --- dengan laju
yang lebih tajam, tetapi hanya bagi distribusi ini.

\textbf{9.} Argumen penukarannya sama sekali tak memakai apa pun dari distribusi
$X_i$-nya di luar $\E X_i = 0$, $\E X_i^2 = 1$ dan
keberhinggaan $\E\abs{X_i}^3$: sehingga seandainya kita mengganti Gaussnya
$N_i$ dengan keluarga i.i.d.\ lain yang dua momen pertamanya
sama dan momen ketiganya berhingga, teleskop yang sama akan
membatasi $\abs{\E f(\text{jumlah}_X) - \E f(\text{jumlah}_Y)}$ oleh
$O(n^{-1/2})$ bagi setiap $f$ yang mulus. Jadi statistik mulus atas
jumlah bebas yang besar bersifat \emph{universal}: yakni sampai
galat yang terkuantifikasi, ia bergantung pada distribusi sukunya hanya
lewat dua bilangan. Inilah asas invariansinya:
buktikanlah sebuah teorema limit bagi distribusi yang paling terhitungkan (yakni
Gaussnya, yang segalanya persis di sana), lalu pindahkanlah ke
setiap distribusi lewat penukaran. Skema yang sama --- dengan jumlahnya diganti
fungsional yang lebih rumit --- menjalankan hukum setengah lingkaran Wigner
bagi matriks acak, keuniversalan akar polinomial
acak, dan sebagian besar peluang modern; sedangkan teorema limit
pusatnya adalah contohnya yang pertama dan paling sederhana.

\textbf{10.} Variabel $Z$ merupakan fungsi Borel $(N_{i+1}, \dots,
N_n)$ saja, sedangkan $A = W_i + \theta h - Z$ dan $h$ merupakan
fungsi variabel sisa keluarga bebasnya
$(X_1, \dots, X_n, N_1, \dots, N_n)$: sehingga menurut
asas koalisinya (\cref{thm:b3:probability:independence}),
$Z$ bebas dari $(A, h)$. Lalu sebagai jumlah
$N_j/\sqrt n \sim \mathcal N(0, \frac1n)$ yang bebas, $Z \sim \mathcal
N(0, s^2)$ dengan $s^2 = \frac{n-i}n$ (\cref{exo:b3:clt:3}),
yang berkepadatan terbatas oleh $\frac1{s\sqrt{2\pi}}$. Sedangkan distribusi
$((A, h), Z)$ adalah hasil kali kedua distribusi marginalnya, sehingga
Tonelli (lewat transfernya) membekukan blok pertamanya: dengan $G(a) =
\E\abs{g(a + Z)} = \int\abs{g(a + z)}\,\varphi_s(z)\,\dd z
\leq \frac{\norm g_{L^1}}{s\sqrt{2\pi}}$ bagi setiap $a$,
\[
\E\bigl[\abs h^3\abs{g(A + Z)}\bigr] =
\E\bigl[\abs h^3\,G(A)\bigr] \leq
\frac{\norm g_{L^1}}{s\sqrt{2\pi}}\,\E\abs h^3
= \sqrt{\frac{n}{2\pi(n-i)}}\,\norm g_{L^1}\,\E\abs h^3 .
\]

\textbf{11.} Bentuk integral rumus Taylornya menyusul
lewat mengintegralkan $f(w + h) - f(w) =
h\int_0^1f'(w + \theta h)\,\dd\theta$ secara parsial dua kali terhadap
$\theta$. Lalu dengan mengambil \omterm{def:b3:probability:space}{nilai harapannya} pada penukaran ke-$i$,
orde $0, 1, 2$-nya meniadakan diri persis seperti pada pertanyaan 3, sedangkan
kedua sisanya (bagi $h = X_i/\sqrt n$
dan $N_i/\sqrt n$) terbatas, bagi $i \leq n - 1$, menurut
pertanyaan 10 dengan $g = f'''$:
\[
\bigl|\E f(H_i) - \E f(H_{i-1})\bigr| \leq
\int_0^1\frac{(1-\theta)^2}2\,\dd\theta\;
\sqrt{\frac{n}{2\pi(n-i)}}\,\norm{f'''}_{L^1}
\frac{\beta + \gamma}{n^{3/2}}
= \frac{\beta + \gamma}{6\,n^{3/2}}
\sqrt{\frac{n}{2\pi(n-i)}}\,\norm{f'''}_{L^1} .
\]
Lalu menjumlahkannya, dengan $\sum_{i=1}^{n-1}\sqrt{\frac n{n-i}} =
\sqrt n\sum_{m=1}^{n-1}m^{-1/2} \leq 2n$, lalu menambahkan
batas pertanyaan 3 bagi penukaran terakhirnya ($i = n$, sebab tak ada \omterm{def:b3:clt:gaussianvector}{Gauss}
yang tersisa):
\[
\bigl|\E f(T_n) - \E f(G_n)\bigr| \leq
\frac{\beta + \gamma}{3\sqrt{2\pi}}\cdot
\frac{\norm{f'''}_{L^1}}{\sqrt n} +
\frac{M_3(\beta + \gamma)}{6\,n^{3/2}} .
\]
Untuk landaiannya: $\psi_\delta'''(x) =
\delta^{-3}\psi'''\bigl(\frac{x - t}\delta\bigr)$, sehingga
$M_3 = K\delta^{-3}$ dengan $K = \norm{\psi'''}_\infty$ dan
$\norm{\psi_\delta'''}_{L^1} = \delta^{-2}
\norm{\psi'''}_{L^1} = K_1\delta^{-2}$ (lewat substitusinya). Lalu
apitan pertanyaan 7 memberikan
\[
\sup_t\,\bigl|\P(T_n \leq t) - \Phi(t)\bigr| \leq
\frac{K_1(\beta + \gamma)}{3\sqrt{2\pi}}\cdot
\frac1{\delta^2\sqrt n} +
\frac{K(\beta + \gamma)}{6}\cdot
\frac1{\delta^3n^{3/2}} + \frac\delta{\sqrt{2\pi}} .
\]
Pada $\delta = n^{-1/6}$ suku pertama dan ketiganya bernilai
$O(n^{-1/6})$ sedangkan yang tengah $O(n^{-1})$: yakni laju
seragam $O(n^{-1/6})$, yang sejati lebih baik daripada $n^{-1/8}$
pertanyaan 7 --- sebab paruh \omterm{def:b3:clt:gaussianvector}{Gauss} hibridanya mengerjakan
pemulusan tambahannya.

\textbf{12.} Berlaku $\E N_1^3 = 0$ (sebab integrannya ganjil), sedangkan
pengintegralan parsial memberikan $\E N_1^4 = 3\,\E N_1^2 = 3$
(sebab $\int x^3\cdot x\varphi(x)\dd x =
3\int x^2\varphi$). Lalu bagi $f$ berkelas $\mathcal C^4$ dengan
turunan yang terbatas, uraikanlah setiap penukaran sampai orde keempat: maka
suku orde ketiganya mengangkut faktor $\E X_i^3 - \E N_i^3 =
0$ (sebab \omterm{def:b3:probability:independence}{kebebasannya} memfaktorkannya seperti pada pertanyaan 3), sehingga hanya
sisa orde keempatnya
$\int_0^1\frac{(1-\theta)^3}6f^{(4)}(w + \theta
h)h^4\dd\theta$ yang bertahan, dengan $\int_0^1
\frac{(1-\theta)^3}6\dd\theta = \frac1{24}$ dan $\E h^4 =
\beta_4n^{-2}$ atau $3n^{-2}$. Lalu pertanyaan 10 (dengan $g =
f^{(4)}$) membatasi penukaran $i \leq n - 1$, dan menjumlahkannya seperti
pada pertanyaan 11:
\[
\bigl|\E f(T_n) - \E f(G_n)\bigr| \leq
\frac{\beta_4 + 3}{12\sqrt{2\pi}}\cdot
\frac{\norm{f^{(4)}}_{L^1}}{n} +
\frac{M_4(\beta_4 + 3)}{24\,n^2} .
\]
Dengan $\norm{\psi_\delta^{(4)}}_{L^1} = K_2\delta^{-3}$ dan
$M_4 = K'\delta^{-4}$, batas fungsi distribusinya
menjadi $C\delta^{-3}n^{-1} + C'\delta^{-4}n^{-2} +
\frac\delta{\sqrt{2\pi}}$; lalu di $\delta = n^{-1/4}$ suku
luarnya bernilai $O(n^{-1/4})$ dan yang tengah $O(n^{-1})$:
jadi lajunya $O(n^{-1/4})$.

\textbf{13.} Dengan $k$ momen yang cocok, sisa yang bertahan
per penukaran berorde $\E\abs h^{k+1} \asymp
n^{-(k+1)/2}$; sedangkan batas \omterm{def:b3:clt:gaussianvector}{Gauss} tersembunyinya membebankan
$\norm{f^{(k+1)}}_{L^1}$ dan jumlah atas penukarannya menyumbang
faktor $2n$, sehingga memberikan $\asymp\norm{f^{(k+1)}}_{L^1}\,
n^{-(k-1)/2}$ bagi $f$ yang mulus. Lalu landaiannya berharga
$\norm{\psi_\delta^{(k+1)}}_{L^1} \asymp \delta^{-k}$, sehingga
galat fungsi distribusinya $\asymp
\delta^{-k}n^{-(k-1)/2} + \delta$, yang berimbang di $\delta =
n^{-(k-1)/(2k+2)}$: jadi lajunya $n^{-(k-1)/(2k+2)}$, yaitu
$n^{-1/6}$ bagi $k = 2$, $n^{-1/4}$ bagi $k = 3$, dan menuju
$n^{-1/2}$ hanya saat $k \to \infty$ --- padahal $k \geq 4$
akan memaksa $\E X_1^4 = 3$ dan seterusnya, yakni distribusi yang
sudah meniru Gaussnya. Kejenuhannya bersifat
struktural: sebab penukarannya menjumlahkan $n$ galat penukaran \emph{dalam
nilai mutlak}, dengan melepaskan semua peniadaan antar
penukarannya. Sedangkan bukti Fouriernya membandingkan \omterm{def:b3:clt:cf}{fungsi karakteristiknya},
yang galatnya muncul beserta fase berayunnya;
lalu ketaksamaan pemulusan Esseen mengubah $\abs{\varphi_{T_n}
- \varphi_N}$, yang diintegralkan terhadap $\frac{\dd\xi}{\abs\xi}$,
menjadi batas fungsi distribusinya dengan ongkos logaritmik
belaka, dan menghasilkan $C\beta n^{-1/2}$ milik Berry--Esseen dari
tiga momen. Jadi penggantiannya menukar keoptimuman dengan
ketegaran --- dan, seperti yang ditunjukkan Bagian VII, dengan kemudahtularannya.

\textbf{14.} Matriks $\Sigma$ setangkup dan semidefinit positif;
lalu dengan $\Sigma = PDP^{\mathsf T}$ (dengan $P$ ortogonal, $D \geq 0$
diagonal, \cref{exo:b3:submanifolds:8}), matriks setangkup $C =
P\sqrt DP^{\mathsf T}$ memenuhi $C^2 = \Sigma$. Lalu bagi sembarang $t
\in \R^2$, $\langle t, CZ\rangle = \langle Ct, Z\rangle =
(Ct)_1Z^1 + (Ct)_2Z^2$ merupakan kombinasi linear variabel
\omterm{def:b3:clt:gaussianvector}{Gauss} yang bebas, jadi \omterm{def:b3:clt:gaussianvector}{Gauss}
(\cref{exo:b3:clt:3}): sehingga $N = CZ$ merupakan \omterm{def:b3:clt:gaussianvector}{vektor Gauss}; dengan
rata-rata $0$ dan kovarians $\E[NN^{\mathsf T}] =
C\,\E[ZZ^{\mathsf T}]\,C^{\mathsf T} = CC^{\mathsf T} =
\Sigma$. Untuk momennya: $\norm N^3 \leq (\abs{N_1} +
\abs{N_2})^3 \leq 4(\abs{N_1}^3 + \abs{N_2}^3)$ (menurut kecembungan
$x^3$ pada $\R_+$), sedangkan setiap koordinatnya merupakan
\omterm{def:b3:clt:gaussianvector}{Gauss} real yang bermomen pada setiap orde
(\cref{exo:b3:product:10}): jadi $\gamma' < \infty$. Akhirnya
setiap $\langle t, G_n\rangle = \frac1{\sqrt
n}\sum_i\langle t, N_i\rangle$ merupakan jumlah ternormalkan
$\mathcal N(0, t^{\mathsf T}\Sigma t)$ yang i.i.d., jadi
persis $\mathcal N(0, t^{\mathsf T}\Sigma t)$: sehingga $G_n$ merupakan
\omterm{def:b3:clt:gaussianvector}{vektor Gauss} berrata-rata $0$ dan berkovarians $\Sigma$, dan
distribusinya $\mathcal N(0, \Sigma)$
(\cref{def:b3:clt:gaussianvector}: sebab distribusinya ditentukan
oleh data itu).

\textbf{15.} Misalkan $\phi(t) = f(w + th)$, dengan $t \in \intcc01$:
maka $\phi$ bersifat $\mathcal C^3$ dengan
\[
\phi'''(t) = \sum_{j,k,l\in\{1,2\}}\partial_{jkl}f(w +
th)\,h_jh_kh_l, \qquad \abs{\phi'''(t)} \leq
M_3\Bigl(\sum_j\abs{h_j}\Bigr)^3 = M_3(\abs{h_1} +
\abs{h_2})^3 .
\]
Lalu Taylor--Lagrange pada orde $3$ bagi $\phi$ antara $0$ dan
$1$ memberikan ketaksamaan pertamanya; sedangkan Cauchy--Schwarz memberikan
$\abs{h_1} + \abs{h_2} \leq \sqrt2\norm h$, dari sanalah
konstanta $\frac{2\sqrt2M_3}6 = \frac{\sqrt2M_3}3$.

\textbf{16.} Definisikanlah $H_i$ dan $W_i$ seperti pada pertanyaan 1, kini
di $\R^2$; sedangkan argumen koalisinya tak berubah. Lalu pada
penukaran ke-$i$, suku orde pertamanya memberikan
$\sum_j\E[\partial_jf(W_i)]\,(\E X_{i,j} - \E N_{i,j})/
\sqrt n = 0$ sedangkan suku orde keduanya memberikan
$\frac1{2n}\sum_{j,k}\E[\partial_{jk}f(W_i)]\,(\Sigma_{jk}
- \Sigma_{jk}) = 0$: yakni rata-rata dan kovariansnya cocok. Lalu pertanyaan
15 membatasi kedua sisanya:
\[
\bigl|\E f(H_i) - \E f(H_{i-1})\bigr| \leq
\frac{\sqrt2M_3}3\cdot\frac{\E\norm{X_i}^3 +
\E\norm{N_i}^3}{n^{3/2}}
= \frac{\sqrt2M_3(\beta' + \gamma')}{3\,n^{3/2}},
\]
lalu dengan menteleskopkannya atas $n$ penukarannya:
\[
\Bigl|\E f\Bigl(\frac{S_n}{\sqrt n}\Bigr) - \E
f(G_n)\Bigr| \leq \frac{\sqrt2\,M_3(\beta' +
\gamma')}{3\sqrt n},
\qquad G_n \sim \mathcal N(0, \Sigma)\ \text{secara persis} .
\]
Untuk kenaikannya: $\E\norm{T_n}^2 = \E\norm{X_1}^2 =
\operatorname{tr}\Sigma$ (sebab suku silangnya lenyap menurut
\omterm{def:b3:probability:independence}{kebebasan} dan pemusatannya), sehingga $\P(\norm{T_n} > A) \leq
\operatorname{tr}\Sigma/A^2$, dan demikian pula bagi $N$: yakni
keketatannya. Lalu diberikan $f$ yang \omterm{def:b3:topology:continuity}{kontinu} terbatas dan $\varepsilon
> 0$, kalikanlah dengan dataran mulus $\chi$ yang sama dengan $1$ pada
bola berjari-jari $A$ dan bertumpu pada jari-jari $A + 1$
(muluskanlah sebuah indikator di $\R^2$,
\cref{thm:b3:lp:regularization}); maka $g = f\chi$ \omterm{def:b3:topology:continuity}{kontinu}
seragam dengan tumpuan \omterm{def:b3:topology:compact}{kompak}, sehingga pemulusan
dua dimensinya $g_\eta$ bersifat $\mathcal C^\infty$ dengan turunan
terbatas pada setiap orde dan $\norm{g - g_\eta}_\infty
\leq \varepsilon$ bagi $\eta$ yang kecil. Lalu rantai tiga-$\varepsilon$
pertanyaan 5 berpindah kata demi kata: $\E f(T_n)
\to \E f(N)$ bagi setiap $f \colon \R^2
\to \R$ yang \omterm{def:b3:topology:continuity}{kontinu} terbatas. Inilah \cref{thm:b3:clt:multiclt} bagi $d = 2$,
yang kini terbukti --- sebab penukarannya mengelakkan teorema
Lévy dua dimensi yang dibiarkan diakui oleh babnya.

\textbf{17.} Bagi $g \colon \R \to \R$ yang \omterm{def:b3:topology:continuity}{kontinu} terbatas,
pemetaan $x \mapsto g(\langle t, x\rangle)$ bersifat \omterm{def:b3:topology:continuity}{kontinu}
terbatas pada $\R^2$, sehingga pertanyaan 16 memberikan $\E g(\langle
t, T_n\rangle) \to \E g(\langle t, N\rangle)$: jadi setiap
proyeksinya konvergen dalam distribusi, dan $\langle t,
N\rangle \sim \mathcal N(0, t^{\mathsf T}\Sigma t)$. (Inilah
arah mudah Cramér--Wold: bahwa konvergensi bersama
mengakibatkan konvergensi setiap peta linearnya.) Penerapannya:
$V_i = (\xi_i, \xi_i^2 - 1)$ merupakan vektor terpusat i.i.d.\
(sebab $\E\xi_1^2 = 1$), dengan entri kovarians $\V(\xi_1) =
1$, $\operatorname{Cov}(\xi_1, \xi_1^2 - 1) = \E\xi_1^3$
dan $\V(\xi_1^2 - 1) = \E\xi_1^4 - 1$; sedangkan momen ketiganya
$\E\norm{V_1}^3 \leq 4\bigl(\E\abs{\xi_1}^3 +
\E\abs{\xi_1^2 - 1}^3\bigr)$ berhingga bila $\xi_1 \in
L^6$. Lalu pertanyaan 16 menghasilkan limit \omterm{def:b3:clt:gaussianvector}{Gauss} bersama
yang ditampilkan itu, dan \cref{thm:b3:clt:gaussianvector}(2): yakni kedua
koordinat limitnya bebas tepat saat
kovarians $\E\xi_1^3$ lenyap --- jadi bagi distribusi setangkup,
rata-rata empiris dan varians empirisnya terurai
secara asimtotik.

\textbf{18.} Tulislah $g(x) - g(\theta) = (g'(\theta) +
\eta(x))(x - \theta)$ dengan $\eta(x) = \frac{g(x) -
g(\theta)}{x - \theta} - g'(\theta)$ bagi $x \neq \theta$
dan $\eta(\theta) = 0$: sebab terdiferensialkan di $\theta$
tepat berarti $\eta(x) \to 0$ saat $x \to \theta$.
\emph{Langkah 1:} $\hat\theta_n \to \theta$ dalam peluang:
sebab bagi $\varepsilon > 0$ dan sembarang $A > 0$, pada akhirnya
$\varepsilon\sqrt n \geq A$, sehingga $\P(\abs{\hat\theta_n -
\theta} > \varepsilon) \leq \P(\abs{\sqrt
n(\hat\theta_n - \theta)} > A) \to \P(\sigma\abs N > A)$
(sebab fungsi distribusinya konvergen pada titik \omterm{def:b3:topology:continuity}{kekontinuan}
$\pm A$), sedangkan ruas kanannya menuju $0$ saat $A \to
\infty$. \emph{Langkah 2:} $\eta(\hat\theta_n) \to 0$ dalam
peluang: sebab diberikan $\varepsilon' > 0$, pilihlah $\delta$ dengan
$\abs\eta \leq \varepsilon'$ pada $\abs{x - \theta} \leq
\delta$; maka $\P(\abs{\eta(\hat\theta_n)} > \varepsilon')
\leq \P(\abs{\hat\theta_n - \theta} > \delta) \to 0$.
\emph{Langkah 3:}
\[
\sqrt n\bigl(g(\hat\theta_n) - g(\theta)\bigr) =
g'(\theta)\,\sqrt n(\hat\theta_n - \theta) +
\eta(\hat\theta_n)\cdot\sqrt n(\hat\theta_n - \theta) .
\]
Lalu menurut aturan hasil kali Slutsky (\cref{exo:b3:clt:8}, dengan
barisan $\eta(\hat\theta_n) \to 0$ dalam peluang dan
$\sqrt n(\hat\theta_n - \theta)$ yang konvergen dalam distribusi),
suku keduanya konvergen dalam distribusi ke $0\cdot\mathcal N(0,
\sigma^2) = 0$, jadi ke $0$ dalam peluang
(\cref{exo:b3:clt:4}(b)); sedangkan yang pertama konvergen dalam distribusi ke
$g'(\theta)\mathcal N(0, \sigma^2)$ (lewat Slutsky lagi, atau lewat
aturan afin \omterm{def:b3:clt:cf}{fungsi karakteristiknya}); lalu aturan jumlah
Slutsky merakit keduanya: jadi limitnya $\mathcal N(0,
g'(\theta)^2\sigma^2)$.

\textbf{19.} (a) Teorema limit pusatnya memberikan $\sqrt n(\bar X_n - \mu)
\Rightarrow \mathcal N(0, \sigma^2)$; lalu metode deltanya
dengan $g(x) = x^2$ dan $g'(\mu) = 2\mu$ memberikan $\sqrt n(\bar
X_n^2 - \mu^2) \Rightarrow \mathcal N(0, 4\mu^2\sigma^2)$
--- yang merosot (berlimit $0$) saat $\mu = 0$. Dalam hal itu
fluktuasinya hidup satu skala di atasnya: $n\bar X_n^2 =
(\sqrt n\,\bar X_n)^2$, sehingga bagi $t > 0$
\[
\P\bigl(n\bar X_n^2 \leq t\bigr) = \P\bigl(-\sqrt t \leq
\sqrt n\,\bar X_n \leq \sqrt t\bigr) \longrightarrow
\Phi\Bigl(\frac{\sqrt t}\sigma\Bigr) -
\Phi\Bigl(-\frac{\sqrt t}\sigma\Bigr) = \P(\sigma^2N^2
\leq t) :
\]
yakni $n\bar X_n^2 \Rightarrow \sigma^2N^2$, kuadrat sebuah
\omterm{def:b3:clt:gaussianvector}{Gauss} (yakni distribusi ``khi-kuadrat'') --- jadi saat turunan
pertamanya mati, suku orde kedua Taylornya mendiktekan
limit yang tak \omterm{def:b3:clt:gaussianvector}{Gauss}. (b) Di sini $\sqrt n(\hat p_n - p)
\Rightarrow \mathcal N(0, p(1 - p))$ dan $g(p) =
\arcsin\sqrt p$ mempunyai $g'(p) = \frac1{2\sqrt{p(1 - p)}}$, sehingga
$g'(p)^2\,p(1 - p) = \frac14$: jadi limitnya $\mathcal
N(0, \frac14)$ bagi setiap $p \in \intoo01$. Pada skala
$\arcsin$-nya, batang galat $95\%$ asimtotiknya adalah $\pm
\frac{0.98}{\sqrt n}$, yang diketahui di muka --- sedangkan pada
\cref{ex:b3:clt:confidence} lebarnya melibatkan
$\sigma = \sqrt{p(1-p)}$ yang tak diketahui, yang harus dikasuskan terburuk oleh $\frac12$
atau ditaksir: jadi transformasinya \emph{menstabilkan}
variansnya.

\textbf{20.} Misalkan $A^* = \{k : \mu(\{k\}) > \nu(\{k\})\}$
dan $\Delta_k = \mu(\{k\}) - \nu(\{k\})$, sehingga
$\sum_k\Delta_k = 0$. Lalu bagi sembarang $A \subseteq \N$: $\mu(A) -
\nu(A) = \sum_{k\in A}\Delta_k \leq \sum_{k\in
A^*}\Delta_k$, dengan kesamaannya di $A = A^*$; dan karena
bagian positif dan negatif $(\Delta_k)$ bermassa
total sama, $\sum_{A^*}\Delta_k =
\frac12\sum_k\abs{\Delta_k}$. Lalu menukar $\mu, \nu$ menangani
tandanya: jadi $d_{\mathrm{TV}}(\mu, \nu) =
\frac12\sum_k\abs{\Delta_k}$. Untuk penggandengannya: bagi sembarang $A$,
\[
\mu(A) - \nu(A) = \E\bigl[\mathbf 1_A(X) - \mathbf
1_A(Y)\bigr] = \E\bigl[(\mathbf 1_A(X) - \mathbf
1_A(Y))\,\mathbf 1_{X\neq Y}\bigr] \leq \P(X \neq Y),
\]
lalu ambillah supremumnya atas $A$.

\textbf{21.} Kedua distribusinya membebani: pada $k = 0$: $1 - p$ lawan
$\eu^{-p}$, dengan $\eu^{-p} > 1 - p$; pada $k = 1$: $p$ lawan
$p\,\eu^{-p} < p$; pada $k \geq 2$: $0$ lawan sisa
Poissonnya $1 - \eu^{-p} - p\eu^{-p} \geq 0$. Karena itu
\[
d_{\mathrm{TV}} = \tfrac12\bigl[(\eu^{-p} - 1 + p) + (p -
p\eu^{-p}) + (1 - \eu^{-p} - p\eu^{-p})\bigr] =
\tfrac12\bigl(2p - 2p\eu^{-p}\bigr) = p(1 - \eu^{-p}),
\]
sedangkan $1 - \eu^{-p} \leq p$ memberikan batasnya $p^2$.

\textbf{22.} Tulislah $H_{i-1} = W_i + X_i$ dan $H_i = W_i +
Y_i$ dengan $W_i = \sum_{j<i}Y_j + \sum_{j>i}X_j$, yang
bebas dari pasangan $(X_i, Y_i)$ (menurut koalisinya). Lalu bagi $A
\subseteq \N$, dengan menyaratkan pada terbilang banyak nilainya
menurut \omterm{def:b3:probability:independence}{kebebasannya},
\[
\P(H_{i-1} \in A) = \sum_{k\geq0}\P(X_i = k)\,\P(W_i + k
\in A),
\]
dan demikian pula bagi $H_i$ dengan $Y_i$. Lalu dengan mengurangkannya, dengan $c_k
= \P(W_i + k \in A) \in \intcc01$ dan $\Delta_k = \P(X_i =
k) - \P(Y_i = k)$ yang berjumlah nol:
\[
\abs{\P(H_{i-1} \in A) - \P(H_i \in A)} =
\Bigl|\sum_k\Delta_k\bigl(c_k - \tfrac12\bigr)\Bigr| \leq
\tfrac12\sum_k\abs{\Delta_k} =
d_{\mathrm{TV}}\bigl(\mathcal B(1, p_i), \mathcal
P(p_i)\bigr) .
\]
Lalu dengan menteleskopkannya dari $H_0 = S$ ke $H_n = \sum_iY_i \sim
\mathcal P(\lambda)$ (\cref{exo:b3:clt:1}, yang diiterasi) dan
memakai pertanyaan 21:
\[
\abs{\P(S \in A) - \P(\mathcal P(\lambda) \in A)} \leq
\sum_{i=1}^np_i\bigl(1 - \eu^{-p_i}\bigr) \leq
\sum_{i=1}^np_i^2
\qquad\text{bagi setiap } A :
\]
yakni ketaksamaan Le Cam. (Sedangkan batas penggandengan pertanyaan 20 memberikan
jalur alternatif: gandengkanlah setiap pasangannya pada satu variabel
seragam sehingga $\P(X_i \neq Y_i) \leq p_i^2$ lalu batasilah
$\P(S \neq \sum Y_i)$; padahal penukarannya tak memerlukan konstruksi
sama sekali.)

\textbf{23.} (a) Dengan $p_i = \frac\lambda n$:
$d_{\mathrm{TV}}(\text{distribusi }S, \mathcal P(\lambda))
\leq \frac{\lambda^2}n$. Ini mempertajam
\cref{exo:b3:clt:5} tiga kali lipat: yakni galat eksplisit pada
setiap $n$ yang berhingga, keseragaman atas semua kejadian $A$ sekaligus
(bukan satu selang setiap kali), dan tanpa menuntut kesamaan
$p_i$-nya --- melainkan hanya $\sum_ip_i^2$ yang kecil, misalnya $\sum p_i^2 \leq
\lambda\max_ip_i$: yakni \emph{banyak kejadian langka, tak satu pun
yang dominan}. (b) Di sini $n = 500$, $p_i = \frac1{500}$,
$\lambda = 1$: sehingga model Poissonnya keliru paling banyak
$500\cdot\frac1{500^2} = 0.002$ pada setiap kejadiannya; khususnya,
dengan mengambil $A = \{0\}$,
\[
\P(\text{tak ada surat tersesat}) = \Bigl(1 -
\frac1{500}\Bigr)^{500},
\qquad
\Bigl|\P(\text{tak ada surat tersesat}) - \eu^{-1}\Bigr| \leq
0.002,
\]
sehingga jawabannya $\eu^{-1} \approx 0.368$ sampai selisih
terjamin $0.002$ (sedangkan selisih sebenarnya sekitar
$4\cdot10^{-4}$). (c) Jadi
soalnya ditutup pada satu metode dengan dua rezim. Saat $n$
sumbangan yang sebanding masing-masing membawa varians
$\frac1n$, mencocokkan \emph{dua} momen terhadap
Gaussnya membuat galat penukarannya $o(\frac1n)$ masing-masing: sehingga jumlahnya
menjadi \omterm{def:b3:clt:gaussianvector}{Gauss} --- dengan Taylor sebagai alat perbandingan lokalnya.
Sedangkan saat $n$ sumbangannya berupa indikator berpeluang
$p_i$, mencocokkan \emph{rata-ratanya} terhadap atom Poissonnya
membuat setiap penukarannya berharga $p_i^2$: sehingga cacahan kejadian langkanya menjadi
Poisson --- dengan variasi total sebagai perbandingan lokal
yang persis. Jadi hibrida yang sama, teleskop yang sama, taksiran lokal
yang berbeda: penggantiannya adalah sebuah strategi, bukan sebuah teorema, dan
limit \omterm{def:b3:clt:gaussianvector}{Gauss} serta Poissonnya adalah dua dividennya yang
tertua.

\textbf{24.} Teorema limit pusatnya memberikan $\sqrt n(\bar X_n - \mu)
\Rightarrow \mathcal N(0, \sigma^2)$, sedangkan $g(x) = \ln x$
terdiferensialkan di $\mu > 0$ dengan $g'(\mu) = \frac1\mu$: sehingga
metode deltanya (Bagian VI) menghasilkan $\sqrt n(\ln\bar X_n - \ln\mu)
\Rightarrow \mathcal N(0, \sigma^2/\mu^2)$. Lalu dengan membuka
selang $\abs{\ln\bar X_n - \ln\mu} \leq
\frac{1.96\,\sigma}{\mu\sqrt n}$ lewat pemangkatan:
\[
\mu \in \bar X_n\cdot
\eu^{\pm1.96\,\sigma/(\mu\sqrt n)}
\qquad\text{dengan peluang asimtotik } 95\%
\]
(dalam praktiknya $\sigma/\mu$ diganti oleh versi
empirisnya, lewat Slutsky seperti pada \cref{exo:b3:clt:8}). Adapun
selang perkaliannya adalah yang alami saat datanya
positif dengan galat yang sebanding dengan \omterm{def:b3:measure:measure}{ukurannya} ---
pendapatan, konsentrasi, waktu paruh: yakni besaran yang hidup
pada skala log, yang selang aditif setangkupnya bahkan bisa
melewati nol di sana.

\textbf{25.} Berlaku $\E[(X - p)^3] = (1-p)^3p + (-p)^3(1 - p) =
p(1-p)\bigl[(1-p)^2 - p^2\bigr] = p(1-p)(1 - 2p)$. Lalu pada
analisis penukarannya (Bagian IV), suku galat yang utama setelah
mencocokkan dua momen mengangkut momen ketiga yang \emph{bertanda}:
sebab bagi $p < \frac12$ ia positif (yakni distribusinya condong
ke kanan: dengan gerak jauh besar yang langka di atas rata-ratanya), sehingga
hampiran normalnya menempatkan massanya secara keliru dan sistematis ---
dengan mengecilkan ekor kiri yang pendek dan membesarkan
yang kanan --- bergalat berorde $n^{-1/2}$; sedangkan di $p =
\frac12$ momen ketiganya lenyap, Bernoullinya cocok
dengan Gaussnya sampai orde ketiga, dan lajunya membaik (yakni pertanyaan
momen-cocok Bagian IV). Secara numerik: $\P(S = 0) =
0.9^{20} = 0.1216$, sedangkan Gaussnya
$\mathcal N(2, 1.8)$ memberikan $\Phi\bigl(\frac{0.5 -
2}{\sqrt{1.8}}\bigr) = \Phi(-1.118) \approx 0.132$: jadi
kurva normalnya, yang tak tahu akan dinding di $0$ dan akan
kemencengan ke kanannya, menaruh terlalu banyak massa di dasarnya --- yakni
tanda galat yang diramalkan, yang sudah terlihat di $n = 20$.

\end{solution}
